
\documentclass[journal]{IEEEtran}

\usepackage{cite}
\usepackage{amsmath,amssymb,amsfonts}
\usepackage{graphicx}
\usepackage{textcomp}
\usepackage{xcolor}
\usepackage{booktabs}
\usepackage{array}
\usepackage{url}
\usepackage{stfloats}

\begin{document}

\title{A Lightweight Balanced Random Forest Intrusion Detection System Using Statistical and Flow-Based Features for Resource-Constrained IoT Networks}

\author{Amiya~Kumar~Sahu%
\thanks{Amiya Kumar Sahu is with the Department of Computer Science and Applications, Sambalpur University, Odisha, India (e-mail: amiyakumar20072@gmail.com).}%
}

\markboth{IEEE-Style Journal Manuscript,~2026}%
{Sahu: Lightweight Balanced Random Forest IDS for Resource-Constrained IoT Networks}

\maketitle

\begin{abstract}
Internet of Things (IoT) networks contain large numbers of heterogeneous and resource-constrained devices, making intrusion detection difficult to deploy at gateways and edge monitoring points. Although machine learning-based intrusion detection systems (IDSs) have reported high accuracy on recent IoT datasets, their practical value depends on how they handle class imbalance, false positives, feature extraction cost, and attack-category complexity. This paper presents a balanced and feature-optimized Random Forest IDS using statistical and flow-based features from the CICIoT2023 dataset. Sixty-three merged CICIoT2023 CSV files were processed using chunk-wise loading, data cleaning, memory optimization, binary label construction, balanced sampling, and feature scaling. A balanced binary dataset of 200,000 records was created with equal benign and attack representation. Random Forest feature importance was used to compare Top-5, Top-10, and Top-15 feature subsets against a full 39-feature baseline, and the selected Top-15 subset was further evaluated with Decision Tree, Extra Trees, Logistic Regression, Gaussian Naive Bayes, and HistGradientBoosting classifiers. The Top-15 Random Forest model achieved 98.85\% accuracy and 98.84\% F1-score, closely matching the full-feature Random Forest model, which achieved 98.86\% accuracy and 98.85\% F1-score. The reduced model lowered the input feature dimension by 61.54\% and generated only 17 false positives among 20,000 benign test samples, corresponding to an approximate false positive rate of 0.085\%. As an extended diagnostic analysis, multiclass attack-category classification was evaluated using capped class-aware sampling. In this setting, the full 39-feature Random Forest achieved 84.85\% accuracy and 73.56\% Macro-F1, while the Top-15 Random Forest achieved 81.38\% accuracy and 70.61\% Macro-F1. The results show that compact flow-level features are highly effective for lightweight binary IDS, whereas detailed attack-family classification benefits from richer feature representation.
\end{abstract}

\begin{IEEEkeywords}
IoT security, intrusion detection system, CICIoT2023, Random Forest, feature selection, class imbalance, lightweight IDS, edge security.
\end{IEEEkeywords}


\section{Introduction}
IoT infrastructures increasingly consist of low-cost, heterogeneous, and often physically exposed devices deployed across health care monitoring, industrial sensing, transportation, smart homes, and smart city services \cite{alfuqaha2015iot}. This deployment model changes the security problem in two important ways. First, large numbers of devices communicate through diverse protocols, gateways, and edge services. Second, many nodes operate with limited memory, processing capacity, and energy supply, making heavyweight monitoring and expensive feature extraction difficult to justify in practical deployments.

The security risk is well established. The Mirai botnet demonstrated how poorly protected embedded and IoT devices can be compromised at scale and coordinated into large distributed denial-of-service attacks \cite{antonakakis2017mirai}. Similar IoT environments are exposed to scanning, spoofing, brute-force attempts, flooding, and malware-driven traffic. Signature-based defenses remain useful, but their main limitation is that they are reactive when traffic behavior or attack implementations change. This has motivated extensive work on machine learning (ML)-based intrusion detection systems (IDSs), which can learn attack patterns from labeled network traffic and support detection beyond fixed signatures \cite{khraisat2019survey,gyamfi2022iotids}.

Despite this progress, three issues remain important for deployable IoT IDS design. The first is feature dimensionality. Many IDS pipelines use full feature sets even when a smaller subset may preserve detection performance. The second is evaluation bias. Aggregate accuracy can be misleading when datasets are imbalanced, because a model may perform well on the dominant class while producing weak benign recall or excessive false alarms. The third is operational cost. A gateway-level IDS should be evaluated not only by classification accuracy but also by feature extraction burden, inference behavior, false positive rate, and model interpretability.

CICIoT2023 is a suitable benchmark for studying these issues because it provides recent IoT traffic with 33 attack types collected from a 105-device topology \cite{neto2023ciciot}. However, its scale and attack-dominant distribution require careful preprocessing. Direct random sampling can preserve majority-class dominance and obscure benign-class performance. Therefore, this work treats balanced sampling as a methodological component of the IDS design rather than as a minor preprocessing step.

This paper proposes a lightweight balanced Random Forest IDS using statistical and flow-level CICIoT2023 features. Random Forest is selected because it performs well on structured tabular traffic data and provides feature importance scores that can be used to construct interpretable reduced-feature models \cite{breiman2001random}. The main research question is whether a compact flow-based feature subset, selected after balanced preprocessing, can preserve near-baseline detection performance while reducing false positives and feature dimensionality.

The main contributions are as follows:
\begin{itemize}
    \item A reproducible preprocessing pipeline is developed for 63 CICIoT2023 merged CSV files using chunk-wise loading, cleaning, memory optimization, label processing, balanced sampling, and scaling.
    \item A balanced binary dataset of 200,000 records is constructed with 100,000 benign and 100,000 attack samples, enabling class-wise evaluation rather than accuracy-only assessment.
    \item Random Forest feature importance is used to identify compact statistical and flow-based features suitable for lightweight IoT IDS design.
    \item Top-5, Top-10, and Top-15 reduced-feature Random Forest models are compared with a 39-feature baseline under the same classifier configuration.
    \item The selected Top-15 feature subset is further evaluated using Decision Tree, Extra Trees, Logistic Regression, Gaussian Naive Bayes, and HistGradientBoosting classifiers.
    \item The Top-15 Random Forest model achieves 98.85\% accuracy and 98.84\% F1-score while reducing the feature space by 61.54\% and producing only 17 false positives on 20,000 benign test records.
    \item An extended multiclass attack-category evaluation is conducted using capped class-aware sampling to examine whether the reduced feature subset remains useful for attack-family diagnosis beyond binary detection.
\end{itemize}

The remainder of this paper is organized as follows. Section II reviews related work. Section III explains the dataset and preprocessing pipeline. Section IV presents the proposed IDS framework. Section V discusses the experimental results. Section VI outlines limitations and future work. Section VII concludes the paper.


\section{Related Work}
IoT intrusion detection research has developed around four connected concerns: realistic datasets, deployable learning models, feature reduction, and reliable evaluation under imbalanced traffic. This section summarizes these areas and positions the present work within them.

\subsection{Datasets for IoT Intrusion Detection}
Benchmark datasets shape IDS research because they determine the traffic distributions, attack classes, and evaluation assumptions available to researchers \cite{khraisat2019survey,thakkar2020datasets}. Earlier datasets such as KDD99 and NSL-KDD were valuable for initial IDS studies, but they do not adequately represent contemporary IoT device behavior, protocol diversity, or recent botnet and flooding attacks. More recent datasets address these limitations. Bot-IoT supports IoT botnet and forensic analytics \cite{koroniotis2019botiot}; N-BaIoT focuses on network-based detection of IoT botnet behavior using deep autoencoders \cite{meidan2018nbaiot}; Kitsune demonstrates online anomaly detection using an ensemble of autoencoders \cite{mirsky2018kitsune}; and ToN-IoT and Edge-IIoTset extend evaluation toward edge and IIoT environments \cite{moustafa2020toniot,ferrag2022edgeiiotset}.

CICIoT2023 is particularly relevant to the present study because it contains 33 attack types organized into seven broad categories and was generated using an IoT topology of 105 devices \cite{neto2023ciciot}. Its flow-level representation allows IDS experiments without packet payload inspection, which is desirable for lightweight and privacy-preserving monitoring. At the same time, its large scale and attack-rich distribution make controlled sampling and class-wise evaluation necessary.

\subsection{ML-Based IDS and Deployment Cost}
A wide range of ML methods has been applied to IDS, including Decision Tree, Random Forest, Support Vector Machine, Naive Bayes, XGBoost, LightGBM, and neural networks \cite{khraisat2019survey,hajiheidari2019iotids}. Deep learning models can improve detection in settings that require sequential or representation learning, but they often require longer training time, larger parameter spaces, and higher computational resources \cite{kwon2019deep}. These costs are important in IoT gateways, where a detector must process traffic with limited delay and without excessive feature extraction.

Random Forest offers a practical compromise for structured flow-based traffic data. It is relatively easy to tune, robust on tabular features, and provides feature importance scores that support model interpretation and feature reduction \cite{breiman2001random}. These properties make it suitable for lightweight and interpretable IDS design.

\subsection{Feature Selection and Flow-Level Detection}
Feature selection reduces dimensionality, removes weak variables, and improves interpretability in ML pipelines \cite{guyon2003feature}. In IDS, the benefit is both practical and statistical: fewer features can reduce extraction, storage, preprocessing, and inference costs. Ozer \emph{et al.} showed that lightweight IDS models can be developed from compact feature subsets in the Bot-IoT setting \cite{ozer2021lightweight}. Recent work has also emphasized that feature reduction and imbalance handling are connected because both affect the reliability of IDS decisions \cite{musthafa2024balanced,ma2025fsllm}.

The type of feature is also important. Payload-based features are expensive to extract and may be less useful when traffic is encrypted. Flow-level features such as packet-length statistics, inter-arrival time, rate, header length, flags, and protocol indicators summarize communication behavior without accessing payload content. This paper follows that direction by using Random Forest importance to select a compact set of statistical and flow-based features.

\subsection{Class Imbalance and False Positives}
Class imbalance remains a persistent problem in IDS evaluation. If one class dominates the training distribution, a model may produce high aggregate accuracy while learning minority behavior poorly \cite{chawla2002smote,zoghi2021unswimbalance}. In intrusion detection, benign traffic wrongly classified as malicious produces false alarms, whereas malicious traffic classified as benign produces missed detections. Both errors have operational consequences, but excessive false positives are especially damaging because they reduce confidence in alerts and increase administrative workload.

Balancing strategies and resampling methods have therefore become important in IDS research \cite{musthafa2024balanced,shanmugam2025imbalance}. This paper applies balanced binary sampling before model training and uses class-wise metrics to evaluate whether benign detection improves. The goal is not only to increase accuracy but also to produce a model whose decisions are more useful in operational IoT monitoring.

\subsection{Research Gap}
The literature shows that realistic datasets, feature selection, Random Forest-based interpretability, and imbalance handling have each been studied separately. A remaining gap is the combined evaluation of these components under a resource-aware IDS objective. Many studies report high accuracy, but fewer examine whether a reduced flow-level feature subset can maintain performance after balancing the data and reducing false alarms. This paper addresses that gap by combining balanced sampling, Random Forest feature ranking, reduced-feature modeling, false-positive analysis, and an extended attack-category evaluation on CICIoT2023.


\section{Dataset and Preprocessing Methodology}
\label{sec:dataset_preprocessing}

\subsection{Dataset Description}
The experiments use CICIoT2023, developed by the Canadian Institute for Cybersecurity at the University of New Brunswick. CICIoT2023 contains traffic from a 105-device IoT topology and includes 33 attack types grouped into seven broad categories: DDoS, DoS, Reconnaissance, Web-based attacks, Brute Force, Spoofing, and Mirai \cite{neto2023ciciot}. The present work uses the dataset's merged CSV flow files rather than raw packet captures. This choice is consistent with the lightweight objective because flow records are easier to process at edge monitors than full packet payloads.

\subsection{Preprocessing Challenges}
Three characteristics of the dataset guided the preprocessing design. First, the merged files are large enough that loading all data into memory is costly on moderate hardware. Second, the original traffic distribution is attack-dominant; uncontrolled random sampling can therefore bias learning toward malicious traffic and obscure benign-class performance. Third, the files may contain duplicate rows, missing values, and infinite numerical values. If left untreated, these issues can distort scaling, feature ranking, and classifier behavior.

\subsection{Chunk-Wise Cleaning and Optimization}
The 63 merged files were read in chunks of 100,000 rows. Let $D=\{D_1,D_2,\ldots,D_n\}$ denote the raw files, where $n=63$. Each file $D_i$ was divided into chunks $C_{ij}$:
\begin{equation}
D_i=\{C_{i1},C_{i2},\ldots,C_{ik}\}.
\end{equation}
Each chunk was cleaned before being written to the processed collection. Column names were standardized by removing extra spaces and replacing internal spaces with underscores. Infinite values were first treated as missing values:
\begin{equation}
x=\begin{cases}
NaN, & \text{if }x=\infty\text{ or }x=-\infty,\\
x, & \text{otherwise}.
\end{cases}
\end{equation}
Missing numerical values were filled using the median of the corresponding feature within the chunk. Median imputation was preferred because attack traffic may produce skewed feature distributions. Duplicate records were removed at both chunk and processed-file levels. To reduce memory usage, floating-point columns were converted from \texttt{float64} to \texttt{float32}, and integer columns were downcast when possible.

\subsection{Label Processing, Sampling, and Scaling}
The original CICIoT2023 labels were mapped to a binary target:
\begin{equation}
y=\begin{cases}
0, & \text{if Label = BENIGN},\\
1, & \text{otherwise}.
\end{cases}
\end{equation}
An attack-category label was also retained for extended multiclass evaluation. For the main binary IDS task, the target distinguishes benign traffic from all attack traffic. For the extended attack-category task, labels were grouped into broader categories: BENIGN, DDOS, DOS, MIRAI, SPOOFING, VULNERABILITY\_SCAN, and Other.

After cleaning, a balanced dataset was created with 100,000 benign and 100,000 attack samples:
\begin{equation}
D_b=D_{benign}^{100000}\cup D_{attack}^{100000}.
\end{equation}
The balanced set was split into 160,000 training and 40,000 testing samples using an 80:20 stratified split. Non-feature columns, namely \texttt{Label}, \texttt{Binary\_Label}, and \texttt{Attack\_Category}, were removed. The remaining 39 numerical features were scaled using StandardScaler:
\begin{equation}
z=\frac{x-\mu}{\sigma},
\end{equation}
where $\mu$ and $\sigma$ were estimated only from the training data and then applied to both the training and test sets.

\begin{table}[htbp]
\caption{Summary of the Preprocessed Dataset}
\label{tab:dataset_summary}
\centering
\begin{tabular}{|l|c|}
\hline
\textbf{Property} & \textbf{Value} \\
\hline
Dataset & CICIoT2023 \\
Merged files processed & 63 \\
Final records & 200,000 \\
Benign records & 100,000 \\
Attack records & 100,000 \\
Input features & 39 \\
Training records & 160,000 \\
Testing records & 40,000 \\
Task & Binary classification \\
Train-test split & 80:20 \\
\hline
\end{tabular}
\end{table}

\begin{figure}[htbp]
\centering
\includegraphics[width=0.47\textwidth]{preprocessing_workflow.png}
\caption{Preprocessing workflow used to produce the balanced CICIoT2023 train-test dataset.}
\label{fig:preprocessing_workflow}
\end{figure}

\section{Proposed Lightweight IDS Framework}
\label{sec:framework}

The proposed framework reduces the feature burden of IoT intrusion detection while preserving balanced class-wise performance. It has four stages: full-feature baseline training, feature importance ranking, reduced subset construction, and final model selection.

\subsection{Random Forest Baseline}

For an input vector $X=\{x_1,x_2,\ldots,x_m\}$, each tree $T_i$ in the Random Forest predicts a class label. The ensemble output is obtained by majority voting:
\begin{equation}
\hat{y}=\mathrm{mode}\{T_1(X),T_2(X),\ldots,T_n(X)\}.
\end{equation}
Here, $\hat{y}=0$ denotes benign traffic and $\hat{y}=1$ denotes attack traffic. All Random Forest models used 100 estimators, maximum depth 20, a fixed random state, and parallel execution. This fixed configuration allowed performance changes to be attributed primarily to the feature subsets rather than to classifier tuning.

\subsection{Feature Ranking and Reduced Subsets}

The full 39-feature model was trained first. Feature importance was then computed from the trained forest based on impurity reduction. Let $F=\{f_1,f_2,\ldots,f_m\}$ be the original feature set and $I(f_j)$ the importance score of feature $f_j$. The ranked list is represented as:
\begin{equation}
I(f_1)\geq I(f_2)\geq \ldots \geq I(f_m).
\end{equation}
Three reduced subsets were created from this ranking: Top-5, Top-10, and Top-15. Each subset was evaluated with the same Random Forest settings used for the baseline.

\subsection{Final Lightweight Configuration}

The strongest reduced configuration used the following 15 features: \texttt{Number}, \texttt{ack\_flag\_number}, \texttt{HTTPS}, \texttt{Tot\_size}, \texttt{Rate}, \texttt{Header\_Length}, \texttt{IAT}, \texttt{AVG}, \texttt{ack\_count}, \texttt{Variance}, \texttt{Std}, \texttt{Max}, \texttt{Tot\_sum}, \texttt{Time\_To\_Live}, and \texttt{Min}. These variables summarize packet count behavior, protocol use, timing, size distribution, and acknowledgement patterns. Since they are flow-based, they do not require payload inspection. The final model was selected by jointly considering accuracy, F1-score, false positives, feature count, training time, inference behavior, and model size.

\section{Experimental Results and Discussion}
\label{sec:results}

\subsection{Experimental Setup}

The experiments were implemented using Python, pandas, NumPy, scikit-learn, and matplotlib. The test set contained 20,000 benign and 20,000 attack records. Accuracy, precision, recall, F1-score, balanced accuracy, Matthews correlation coefficient (MCC), false positive rate (FPR), false negative rate (FNR), confusion matrix, training time, inference time, and serialized model size were recorded. Accuracy alone was not treated as sufficient because an IDS can be accurate overall while still producing many false alerts or missing minority-class behavior. In addition to the full and reduced Random Forest experiments, the selected Top-15 feature subset was evaluated with Decision Tree, Extra Trees, Logistic Regression, Gaussian Naive Bayes, and HistGradientBoosting classifiers under the same train-test split. A separate extended multiclass experiment was then performed to examine attack-category classification using a capped sampling strategy.

The metrics are defined as follows:
\begin{equation}
Accuracy=\frac{TP+TN}{TP+TN+FP+FN},
\end{equation}
\begin{equation}
Precision=\frac{TP}{TP+FP},
\end{equation}
\begin{equation}
Recall=\frac{TP}{TP+FN},
\end{equation}
\begin{equation}
F1=\frac{2\times Precision\times Recall}{Precision+Recall}.
\end{equation}

\subsection{Full and Reduced Feature Performance}

Table~\ref{tab:reduced_feature_performance} compares the full-feature model with the three reduced-feature models. The 39-feature Random Forest reached 98.86\% accuracy and 98.85\% F1-score. Reducing the model to five features lowered the training time but also reduced accuracy. The Top-10 model recovered part of this loss, while the Top-15 model almost matched the full baseline.

\begin{table}[htbp]
\caption{Performance Comparison of Full and Reduced Feature Random Forest Models}
\label{tab:reduced_feature_performance}
\centering
\begin{tabular}{|l|c|c|c|}
\hline
\textbf{Model} & \textbf{Accuracy} & \textbf{F1-score} & \textbf{Training Time} \\
\hline
RF Full (39 Features) & 98.86\% & 98.85\% & 12.32 s \\
RF Top-5 & 98.45\% & 98.43\% & 7.02 s \\
RF Top-10 & 98.56\% & 98.54\% & 11.39 s \\
RF Top-15 & 98.85\% & 98.84\% & 11.46 s \\
\hline
\end{tabular}
\end{table}

The Top-15 model reduced the feature dimension from 39 to 15:
\begin{equation}
Feature\ Reduction=\frac{39-15}{39}\times100=61.54\%.
\end{equation}
This result is important because the reduction caused almost no loss in accuracy or F1-score. It indicates that the original feature set contains redundant or weakly useful variables for binary detection.

\subsection{Comparison with Other Classifiers}

To verify that the selected Top-15 feature subset is not effective only for Random Forest, additional classifiers were trained and evaluated on the same balanced CICIoT2023 train-test split. The compared models were Decision Tree, Extra Trees, Logistic Regression, Gaussian Naive Bayes, and HistGradientBoosting. The full 39-feature Random Forest was retained as a reference baseline. This comparison makes the evaluation more useful for journal review because it tests whether the proposed reduced-feature design remains competitive against both simple and ensemble classifiers.

In addition to accuracy and F1-score, the comparison includes balanced accuracy, MCC, ROC-AUC, false positive rate (FPR), and false negative rate (FNR), as shown in Table~\ref{tab:extended_classifier_comparison}. These metrics provide a more complete view of IDS behavior. Balanced accuracy gives equal importance to benign and attack recall. MCC is useful for binary classification because it considers all four confusion-matrix outcomes. ROC-AUC measures the ability of a classifier to separate benign and malicious traffic across decision thresholds.

\begin{table*}[htbp]
\caption{Classifier Comparison with Extended Performance Metrics}
\label{tab:extended_classifier_comparison}
\centering
\scriptsize
\begin{tabular}{lcccccccc}
\toprule
\textbf{Model} & \textbf{Feat.} & \textbf{Acc.} & \textbf{F1} & \textbf{Bal. Acc.} & \textbf{MCC} & \textbf{ROC-AUC} & \textbf{FPR} & \textbf{FNR} \\
\midrule
RF Full-39 & 39 & 98.86 & 98.85 & 0.9886 & 0.9775 & 0.9968 & 0.085 & 2.190 \\
RF Top-15 & 15 & 98.85 & 98.84 & 0.9885 & 0.9772 & 0.9963 & 0.085 & 2.215 \\
HistGB Top-15 & 15 & 98.85 & 98.84 & 0.9885 & 0.9772 & 0.9964 & 0.100 & 2.205 \\
Extra Trees Top-15 & 15 & 98.62 & 98.60 & 0.9862 & 0.9727 & 0.9961 & 0.055 & 2.710 \\
Decision Tree Top-15 & 15 & 98.61 & 98.60 & 0.9861 & 0.9722 & 0.9884 & 0.725 & 2.065 \\
Logistic Reg. Top-15 & 15 & 98.18 & 98.14 & 0.9818 & 0.9641 & 0.9917 & 0.015 & 3.635 \\
Gaussian NB Top-15 & 15 & 97.56 & 97.54 & 0.9756 & 0.9513 & 0.9900 & 1.600 & 3.285 \\
\bottomrule
\end{tabular}
\end{table*}

Accuracy, F1-score, FPR, and FNR in Table~\ref{tab:extended_classifier_comparison} are reported as percentages. Balanced accuracy, MCC, and ROC-AUC are reported as normalized scores. The Random Forest Top-15 model achieved 0.9885 balanced accuracy, 0.9772 MCC, and 0.9963 ROC-AUC. These values were nearly identical to the full 39-feature Random Forest baseline, which achieved 0.9886 balanced accuracy, 0.9775 MCC, and 0.9968 ROC-AUC. This confirms that the Top-15 model preserved almost the same discriminative capability as the full model while reducing the input feature dimension by 61.54\%.

HistGradientBoosting produced a very close result, with 98.85\% accuracy, 98.84\% F1-score, and 0.9964 ROC-AUC, but its false positive rate was slightly higher than that of the Random Forest Top-15 model. Extra Trees produced the lowest FPR, but it also had a higher FNR, indicating more missed attacks. Decision Tree was computationally fast but produced a noticeably higher FPR. Logistic Regression had a very low FPR but a higher FNR and substantially longer training time. Gaussian Naive Bayes was the fastest model, but it had the weakest overall classification performance. Therefore, the proposed Random Forest Top-15 model provides the most balanced tradeoff among detection accuracy, false-positive control, attack recall, interpretability, and feature efficiency.

\subsection{Extended Multiclass Attack Category Evaluation}

Although the main objective of this study is lightweight binary intrusion detection, an extended multiclass experiment was conducted to examine whether the proposed feature-optimized approach can also support attack-family identification. The original attack labels were grouped into seven broader categories: BENIGN, DDOS, DOS, MIRAI, SPOOFING, VULNERABILITY\_SCAN, and Other. Strict equal balancing was avoided because the smallest merged categories contained only around one thousand records, which would discard a large portion of useful traffic from larger classes. Instead, a capped sampling strategy was used: up to 20,000 records were retained from large classes, while all available records were retained from minority classes.

The resulting multiclass dataset contained 70,479 records, consisting of 20,000 BENIGN, 20,000 DDOS, 19,238 DOS, 7,693 MIRAI, 1,346 SPOOFING, 1,108 VULNERABILITY\_SCAN, and 1,094 Other records. The dataset was divided using an 80:20 stratified split, producing 56,383 training records and 14,096 testing records. The same classifier families used in the binary comparison were evaluated for the multiclass task. Macro-F1 was emphasized because it gives equal importance to each class and is therefore more informative when minority attack categories are present.

\begin{table*}[htbp]
\caption{Improved Multiclass Attack Category Classification Results}
\label{tab:multiclass_results}
\centering
\scriptsize
\begin{tabular}{lccccccc}
\toprule
\textbf{Model} & \textbf{Feat.} & \textbf{Acc.} & \textbf{Macro-F1} & \textbf{Weighted-F1} & \textbf{Bal. Acc.} & \textbf{MCC} & \textbf{ROC-AUC} \\
\midrule
RF Full-39 & 39 & 84.85 & 73.56 & 84.25 & 0.6988 & 0.8024 & 0.9720 \\
RF Top-15 & 15 & 81.38 & 70.61 & 80.60 & 0.6723 & 0.7580 & 0.9628 \\
HistGB Top-15 & 15 & 81.48 & 70.18 & 80.72 & 0.6738 & 0.7582 & 0.9662 \\
Extra Trees Top-15 & 15 & 79.80 & 69.52 & 79.56 & 0.7022 & 0.7423 & 0.9604 \\
Decision Tree Top-15 & 15 & 76.61 & 65.52 & 77.31 & 0.6940 & 0.7001 & 0.8708 \\
Logistic Reg. Top-15 & 15 & 65.18 & 51.89 & 67.85 & 0.6118 & 0.5648 & 0.9353 \\
Gaussian NB Top-15 & 15 & 53.23 & 40.29 & 51.75 & 0.5623 & 0.4656 & 0.9237 \\
\bottomrule
\end{tabular}
\end{table*}

Accuracy, Macro-F1, and Weighted-F1 in Table~\ref{tab:multiclass_results} are reported as percentage values, while balanced accuracy, MCC, and ROC-AUC are reported as normalized scores. The full 39-feature Random Forest achieved the strongest multiclass result, with 84.85\% accuracy, 73.56\% Macro-F1, 84.25\% Weighted-F1, 0.8024 MCC, and 0.9720 ROC-AUC. The Top-15 Random Forest achieved 81.38\% accuracy and 70.61\% Macro-F1. The performance reduction relative to the full model indicates that multiclass attack-family classification requires more detailed feature information than binary attack detection.

The difference between Weighted-F1 and Macro-F1 also shows that larger classes such as BENIGN, DDOS, and DOS were classified more reliably than smaller categories such as SPOOFING, VULNERABILITY\_SCAN, and Other. This result is expected because the capped sampling strategy preserves more information from common classes while retaining all available records from rare categories. Therefore, the multiclass experiment should be interpreted as an extended diagnostic analysis rather than as the main lightweight IDS contribution. It shows that the Top-15 feature set remains competitive under feature reduction, but full-feature or class-specific models may be more suitable when the objective is detailed attack-family classification.

\subsection{Comparison with Contemporary Studies}

The proposed approach was also compared with recent IoT IDS studies that emphasize CICIoT2023, feature reduction, class balancing, Random Forest-based learning, or lightweight deployment. Direct numerical comparison across studies must be interpreted carefully because prior works use different train-test splits, sampling strategies, binary or multiclass settings, feature sets, and evaluation metrics. Nevertheless, the comparison in Table~\ref{tab:related_study_comparison} helps position the contribution of the present work.

\begin{table*}[htbp]
\caption{Comparison with Contemporary IoT IDS Studies}
\label{tab:related_study_comparison}
\centering
\scriptsize
\begin{tabular}{p{2.6cm} p{2.0cm} p{3.3cm} p{6.6cm}}
\toprule
\textbf{Study} & \textbf{Dataset} & \textbf{Method} & \textbf{Reported Focus / Result} \\
\midrule
Neto et al. \cite{neto2023ciciot} & CICIoT2023 & Benchmark ML/DL evaluation & Introduced CICIoT2023 with 33 attacks, seven attack classes, and an IoT topology of 105 devices for IoT IDS benchmarking. \\
Wardana et al. \cite{wardana2024frf} & CICIoT2023 & Federated Random Forest with feature selection & Reported approximately 99.68\% accuracy, precision, recall, and F1-score for collaborative IoT intrusion detection. \\
Musthafa et al. \cite{musthafa2024balanced} & UNSW-NB15 / NSL-KDD & Class balancing, feature selection, and ensemble ML & Demonstrated that balancing and feature selection can improve IDS reliability and preprocessing quality. \\
Kouassi et al. \cite{kouassi2025topk} & IoT IDS datasets & Top-K feature selection with RF, XGBoost, and LightGBM & Reported a strong accuracy-efficiency tradeoff using reduced Top-K feature subsets. \\
Proposed work & CICIoT2023 & Balanced RF with Top-15 flow features & Achieved 98.85\% binary accuracy, 98.84\% F1-score, 0.9885 balanced accuracy, 0.9772 MCC, 0.9963 ROC-AUC, and 0.085\% FPR with 61.54\% feature reduction; extended multiclass analysis reached 84.85\% accuracy using RF Full-39 and 81.38\% using RF Top-15. \\
\bottomrule
\end{tabular}
\end{table*}

Neto et al. introduced CICIoT2023 as a modern IoT IDS benchmark and established its relevance for evaluating malicious and benign IoT traffic detection. Wardana et al. proposed a Federated Random Forest model with feature selection for collaborative IoT intrusion detection and reported very high performance, but their main emphasis was privacy-preserving collaboration. Musthafa et al. focused on the joint role of class balancing and feature selection, which supports the methodological direction adopted in the present work. Kouassi et al. further demonstrated that Top-K feature selection can improve the tradeoff between accuracy and computational efficiency in IoT IDS settings.

Compared with these studies, the present work does not claim to replace federated or hybrid ensemble systems. Its contribution is narrower and deployment-oriented: it shows that a class-balanced Random Forest model using only 15 statistical and flow-based features can closely match a full 39-feature baseline, outperform several reduced-feature classifier alternatives in overall F1-score, and maintain a very low FPR. This makes the proposed approach suitable for IoT gateway-level monitoring where interpretability, limited feature extraction, and false-positive control are important.

\subsection{Class-Wise Evaluation}

The Top-15 model was selected as the final lightweight configuration. Its class-wise results are shown in Table~\ref{tab:classwise_top15}. Benign recall reached 1.00, while attack recall remained 0.98. This balance is useful for operational IDS deployment because benign recall directly affects the false alarm burden.

\begin{table}[htbp]
\caption{Class-Wise Performance of the Top-15 Random Forest Model}
\label{tab:classwise_top15}
\centering
\begin{tabular}{|c|c|c|c|c|}
\hline
\textbf{Class} & \textbf{Precision} & \textbf{Recall} & \textbf{F1-score} & \textbf{Support} \\
\hline
Benign (0) & 0.98 & 1.00 & 0.99 & 20000 \\
Attack (1) & 1.00 & 0.98 & 0.99 & 20000 \\
\hline
Accuracy & \multicolumn{4}{c|}{0.99} \\
Macro Avg. & 0.99 & 0.99 & 0.99 & 40000 \\
Weighted Avg. & 0.99 & 0.99 & 0.99 & 40000 \\
\hline
\end{tabular}
\end{table}

\subsection{Confusion Matrix and Error Rates}

The confusion matrix in Table~\ref{tab:confusion_top15} provides a clearer view of the IDS errors. The model correctly classified 19,983 of 20,000 benign records and incorrectly raised only 17 false alarms. It detected 19,557 of 20,000 attack records, while 443 attack records were missed.

\begin{table}[htbp]
\caption{Confusion Matrix of the Top-15 Random Forest Model}
\label{tab:confusion_top15}
\centering
\begin{tabular}{|c|c|c|}
\hline
\textbf{Actual / Predicted} & \textbf{Benign} & \textbf{Attack} \\
\hline
Benign & 19983 & 17 \\
Attack & 443 & 19557 \\
\hline
\end{tabular}
\end{table}

The false positive rate and false negative rate are:
\begin{equation}
FPR=\frac{FP}{FP+TN}=\frac{17}{17+19983}=0.00085,
\end{equation}
\begin{equation}
FNR=\frac{FN}{FN+TP}=\frac{443}{443+19557}=0.02215.
\end{equation}
Thus, the selected model produced an approximate FPR of 0.085\% and an FNR of 2.215\%. The low FPR is the more operationally important finding because false alarms are a major reason IDS systems become difficult to use in practice.

\subsection{Effect of Balanced Sampling}

The earlier imbalanced experiment produced high attack recall but weaker benign recall. Table~\ref{tab:imbalance_effect} compares that behavior with the balanced Top-15 model. The full imbalanced model reached only 0.81 benign recall, and the imbalanced Top-10 model reached 0.79. After balanced sampling, benign recall increased to 1.00 and false positives decreased to 17. Attack recall decreased slightly from 0.99 to 0.98, but the improvement in false alarm behavior is substantial.

\begin{table}[htbp]
\caption{Effect of Balanced Sampling on Class-Wise Detection}
\label{tab:imbalance_effect}
\centering
\begin{tabular}{|l|c|c|c|}
\hline
\textbf{Experiment} & \textbf{Benign Recall} & \textbf{Attack Recall} & \textbf{False Positives} \\
\hline
Imbalanced RF Full & 0.81 & 0.99 & 234 \\
Imbalanced RF Top-10 & 0.79 & 0.99 & 251 \\
Balanced RF Top-15 & 1.00 & 0.98 & 17 \\
\hline
\end{tabular}
\end{table}

\subsection{Computational Observations}

The Top-15 Random Forest model trained in 9.21 seconds in the classifier comparison experiment, compared with 9.45 seconds for the full 39-feature Random Forest. More importantly, inference time decreased from 0.2050 seconds for the full model to 0.1170 seconds for the Top-15 model on the same test set. The reduced model also lowers the number of traffic features that must be extracted and processed at runtime. The serialized model size of the Top-15 Random Forest was 15.98 MB, whereas the full Random Forest was 14.45 MB. This shows that tree-based model size is influenced by tree structure, depth, and node splits rather than feature count alone. For this reason, lightweight IDS should be judged through several indicators: feature count, feature extraction effort, balanced accuracy, MCC, ROC-AUC, class-wise reliability, inference behavior, false positive rate, and not only file size.

\subsection{Discussion}

The experiments show that balanced sampling changes the interpretation of IDS performance. Under the imbalanced setting, the model appeared strong because attack recall was high, but benign detection was weaker. Once the dataset was balanced, the model's false positive behavior improved sharply. This supports the use of class-wise metrics rather than accuracy alone.

The feature results also show that the CICIoT2023 flow feature space contains redundancy for binary detection. A 15-feature model delivered almost the same performance as the full 39-feature model. The selected variables mainly describe timing, size, rate, protocol, and acknowledgement behavior. Since these properties can be obtained from flow records, the proposed approach is more practical for edge monitoring than approaches that depend on payload inspection. The classifier comparison further strengthens this claim: among the reduced-feature alternatives, Random Forest Top-15 achieved the best overall F1-score while maintaining the same FPR as the full Random Forest baseline. The extended metrics reinforce the same conclusion, as RF Top-15 obtained 0.9885 balanced accuracy, 0.9772 MCC, and 0.9963 ROC-AUC. The comparison with contemporary studies also shows that this work contributes a compact and interpretable alternative to more complex federated, stacking, or hybrid architectures. The extended multiclass experiment adds a second interpretation: binary intrusion detection and attack-family diagnosis have different feature requirements. The Top-15 subset is highly effective for lightweight binary detection, but the full 39-feature representation performs better when the classifier must separate multiple attack categories. Overall, the Top-15 Random Forest model provides the best balance between detection quality, false-positive control, interpretability, and resource-awareness for the binary IDS objective, while multiclass categorization remains a more demanding diagnostic extension.

\section{Limitations and Future Work}
\label{sec:limitations}

The present study has four limitations. First, the main contribution is optimized for binary intrusion detection; the multiclass experiment was included as an extended diagnostic analysis and showed lower Macro-F1 than the binary setting. Second, minority attack categories such as SPOOFING, VULNERABILITY\_SCAN, and Other remain relatively small even after capped sampling, which affects macro-level performance. Third, the evaluation uses a static train-test split, whereas operational IoT traffic may drift over time. Fourth, the model was evaluated on a workstation pipeline rather than deployed on physical gateway hardware.

Future work will strengthen multiclass attack-family classification using improved minority-class handling, class-specific feature analysis, and possibly hierarchical detection in which a lightweight binary detector is followed by a separate attack-category classifier. Concept drift methods such as ADWIN, DDM, and EDDM will also be investigated to monitor changes in traffic distribution. Online and incremental learning will be explored so that the IDS can adapt without complete retraining. In addition, federated learning will be studied for privacy-preserving collaboration among IoT gateways. Finally, the Top-15 feature subset will be validated on additional datasets such as Bot-IoT, N-BaIoT, ToN-IoT, and Edge-IIoTset \cite{koroniotis2019botiot,meidan2018nbaiot,moustafa2020toniot,ferrag2022edgeiiotset}, followed by real-time testing on gateway-class hardware such as Raspberry Pi.

\section{Conclusion}
\label{sec:conclusion}

This paper presented a balanced and feature-optimized Random Forest IDS for resource-constrained IoT networks. Using CICIoT2023, 63 merged CSV files were processed into a balanced binary dataset of 200,000 records. The full-feature Random Forest baseline achieved 98.86\% accuracy and 98.85\% F1-score. Feature ranking then showed that a much smaller subset could preserve nearly the same detection capability.

Among the reduced models, the Top-15 feature configuration achieved 98.85\% accuracy, 98.84\% F1-score, 0.9885 balanced accuracy, 0.9772 MCC, and 0.9963 ROC-AUC while reducing the input dimension by 61.54\%. Its confusion matrix showed only 17 false positives among 20,000 benign test records, corresponding to an approximate FPR of 0.085\%. Additional comparison with Decision Tree, Extra Trees, Logistic Regression, Gaussian Naive Bayes, and HistGradientBoosting showed that the Random Forest Top-15 model obtained the best F1-score among the reduced-feature classifiers while preserving the low false-positive behavior of the full Random Forest baseline. The contemporary study comparison further positions the work as a compact and interpretable alternative for balanced IoT IDS evaluation. An extended multiclass experiment showed that RF Full-39 achieved 84.85\% accuracy and 73.56\% Macro-F1, while RF Top-15 achieved 81.38\% accuracy and 70.61\% Macro-F1. This indicates that compact features are sufficient for lightweight binary detection, whereas detailed attack-family classification benefits from richer feature representation. Overall, these results show that balanced sampling and flow-level feature reduction can jointly improve the practical value of IoT IDS models. The proposed approach provides an interpretable and resource-aware basis for gateway-level intrusion detection without relying on packet payload inspection.


\section*{Data and Code Availability}
The CICIoT2023 dataset used in this study is publicly available from the Canadian Institute for Cybersecurity, University of New Brunswick. Because of dataset size and redistribution considerations, the raw CICIoT2023 files are not redistributed with this manuscript. The preprocessing scripts, model training code, feature-selection workflow, result tables, and figures will be made available in a public GitHub repository upon submission or acceptance. Users should download the dataset from the official source and place the merged CSV files in the repository's designated raw-data directory before running the scripts.

\bibliographystyle{IEEEtran}
\bibliography{references_final_verified}

\end{document}
